\section{Relative glueing}
\label{section-relative-glueing}

\noindent
The following lemma is relevant in case we are trying to construct a
scheme $X$ over $S$, and we already know how to construct the restriction
of $X$ to the affine opens of $S$. The actual result is completely general
and works in the setting of (locally) ringed spaces, although our proof
is written in the language of schemes.

\begin{lemma}
\label{lemma-relative-glueing}
Let $S$ be a scheme.
Let $\mathcal{B}$ be a basis for the topology of $S$.
Suppose given the following data:
\begin{enumerate}
\item For every $U \in \mathcal{B}$ a scheme $f_U : X_U \to U$ over $U$.
\item For $U, V \in \mathcal{B}$ with $V \subset U$ a morphism
$\rho^U_V : X_V \to X_U$ over $U$.
\end{enumerate}
Assume that
\begin{enumerate}
\item[(a)] each $\rho^U_V$ induces an isomorphism
$X_V \to f_U^{-1}(V)$ of schemes over $V$,
\item[(b)] whenever $W, V, U \in \mathcal{B}$, with
$W \subset V \subset U$ we have $\rho^U_W = \rho^U_V \circ \rho ^V_W$.
\end{enumerate}
Then there exists a morphism $f : X \to S$ of schemes
and isomorphisms $i_U : f^{-1}(U) \to X_U$ over $U \in \mathcal{B}$
such that for $V, U \in \mathcal{B}$ with $V \subset U$ the composition
$$
\xymatrix{
X_V \ar[r]^{i_V^{-1}} &
f^{-1}(V) \ar[rr]^{inclusion} & &
f^{-1}(U) \ar[r]^{i_U} &
X_U
}
$$
is the morphism $\rho^U_V$. Moreover $X$ is unique up to
unique isomorphism over $S$.
\end{lemma}

\begin{proof}
To prove this we will use Schemes, Lemma \ref{schemes-lemma-glue-functors}.
First we define a contravariant functor $F$ from the category of schemes
to the category of sets. Namely, for a scheme $T$ we set
$$
F(T) =
\left\{
\begin{matrix}
(g, \{h_U\}_{U \in \mathcal{B}}),
\ g : T \to S, \ h_U : g^{-1}(U) \to X_U, \\
f_U \circ h_U = g|_{g^{-1}(U)},
\ h_U|_{g^{-1}(V)} = \rho^U_V \circ h_V
\ \forall\ V, U \in \mathcal{B}, V \subset U
\end{matrix}
\right\}.
$$
The restriction mapping $F(T) \to F(T')$ given a morphism
$T' \to T$ is just gotten by composition.
For any $W \in \mathcal{B}$ we consider the subfunctor
$F_W \subset F$ consisting of those systems $(g, \{h_U\})$
such that $g(T) \subset W$.

\medskip\noindent
First we show $F$ satisfies the sheaf property for the Zariski topology.
Suppose that $T$ is a scheme, $T = \bigcup V_i$ is an open covering,
and $\xi_i \in F(V_i)$ is an element such that
$\xi_i|_{V_i \cap V_j} = \xi_j|_{V_i \cap V_j}$.
Say $\xi_i = (g_i, \{h_{i, U}\})$. Then we immediately see that
the morphisms $g_i$ glue to a unique global morphism
$g : T \to S$. Moreover, it is clear that
$g^{-1}(U) = \bigcup g_i^{-1}(U)$. Hence the morphisms
$h_{i, U} : g_i^{-1}(U) \to X_U$ glue to a unique morphism
$h_U : g^{-1}(U) \to X_U$. It is easy to verify that the system
$(g, \{h_U\})$ is an element of $F(T)$. Hence $F$ satisfies the
sheaf property for the Zariski topology.

\medskip\noindent
Next we verify that each $F_W$, $W \in \mathcal{B}$ is representable.
Namely, we claim that the transformation of functors
$$
F_W \longrightarrow \Mor(-, X_W), \ (g, \{h_U\}) \longmapsto h_W
$$
is an isomorphism. To see this suppose that $T$ is a scheme and
$\alpha : T \to X_W$ is a morphism. Set $g = f_W \circ \alpha$.
For any $U \in \mathcal{B}$ such that $U \subset W$ we can
define $h_U : g^{-1}(U) \to X_U$ be the composition
$(\rho^W_U)^{-1} \circ \alpha|_{g^{-1}(U)}$. This works because
the image $\alpha(g^{-1}(U))$ is contained in $f_W^{-1}(U)$ and
condition (a) of the lemma. It is clear that
$f_U \circ h_U = g|_{g^{-1}(U)}$ for such a $U$.
Moreover, if also $V \in \mathcal{B}$ and $V \subset U \subset W$,
then $\rho^U_V \circ h_V = h_U|_{g^{-1}(V)}$ by property (b)
of the lemma. We still have to define $h_U$ for an arbitrary
element $U \in \mathcal{B}$. Since $\mathcal{B}$ is a basis for
the topology on $S$ we can find an open covering
$U \cap W = \bigcup U_i$ with $U_i \in \mathcal{B}$. Since $g$ maps into $W$
we have
$g^{-1}(U) = g^{-1}(U \cap W) = \bigcup g^{-1}(U_i)$.
Consider the morphisms
$h_i = \rho^U_{U_i} \circ h_{U_i} : g^{-1}(U_i) \to X_U$.
It is a simple matter to use condition (b) of the lemma
to prove that
$h_i|_{g^{-1}(U_i) \cap g^{-1}(U_j)} = h_j|_{g^{-1}(U_i) \cap g^{-1}(U_j)}$.
Hence these morphisms glue to give the desired morphism
$h_U : g^{-1}(U) \to X_U$. We omit the (easy) verification that
the system $(g, \{h_U\})$ is an element of $F_W(T)$ which
maps to $\alpha$ under the displayed arrow above.

\medskip\noindent
Next, we verify each $F_W \subset F$ is representable by open immersions.
This is clear from the definitions.

\medskip\noindent
Finally we have to verify
the collection $(F_W)_{W \in \mathcal{B}}$ covers $F$.
This is clear by construction and the fact that $\mathcal{B}$ is
a basis for the topology of $S$.

\medskip\noindent
Let $X$ be a scheme representing the functor $F$.
Let $(f, \{i_U\}) \in F(X)$ be a ``universal family''.
Since each $F_W$ is representable by $X_W$ (via the morphism of functors
displayed above) we see that $i_W : f^{-1}(W) \to X_W$
is an isomorphism as desired. The lemma is proved.
\end{proof}

\begin{lemma}
\label{lemma-relative-glueing-sheaves}
Let $S$ be a scheme.
Let $\mathcal{B}$ be a basis for the topology of $S$.
Suppose given the following data:
\begin{enumerate}
\item For every $U \in \mathcal{B}$ a scheme $f_U : X_U \to U$ over $U$.
\item For every $U \in \mathcal{B}$ a quasi-coherent sheaf $\mathcal{F}_U$
over $X_U$.
\item For every pair $U, V \in \mathcal{B}$ such that
$V \subset U$ a morphism $\rho^U_V : X_V \to X_U$.
\item  For every pair $U, V \in \mathcal{B}$ such that
$V \subset U$ a morphism
$\theta^U_V : (\rho^U_V)^*\mathcal{F}_U \to \mathcal{F}_V$.
\end{enumerate}
Assume that
\begin{enumerate}
\item[(a)] each $\rho^U_V$ induces an isomorphism
$X_V \to f_U^{-1}(V)$ of schemes over $V$,
\item[(b)] each $\theta^U_V$ is an isomorphism,
\item[(c)] whenever $W, V, U \in \mathcal{B}$, with
$W \subset V \subset U$ we have $\rho^U_W = \rho^U_V \circ \rho ^V_W$,
\item[(d)] whenever $W, V, U \in \mathcal{B}$, with
$W \subset V \subset U$ we have
$\theta^U_W = \theta^V_W \circ (\rho^V_W)^*\theta^U_V$.
\end{enumerate}
Then there exists a morphism of schemes $f : X \to S$
together with a quasi-coherent sheaf $\mathcal{F}$ on $X$
and isomorphisms $i_U : f^{-1}(U) \to X_U$ and
$\theta_U : i_U^*\mathcal{F}_U \to \mathcal{F}|_{f^{-1}(U)}$
over $U \in \mathcal{B}$ such that
for $V, U \in \mathcal{B}$ with $V \subset U$ the composition
$$
\xymatrix{
X_V \ar[r]^{i_V^{-1}} &
f^{-1}(V) \ar[rr]^{inclusion} & &
f^{-1}(U) \ar[r]^{i_U} &
X_U
}
$$
is the morphism $\rho^U_V$, and the composition
\begin{equation}
\label{equation-glue}
(\rho^U_V)^*\mathcal{F}_U
=
(i_V^{-1})^*((i_U^*\mathcal{F}_U)|_{f^{-1}(V)})
\xrightarrow{\theta_U|_{f^{-1}(V)}}
(i_V^{-1})^*(\mathcal{F}|_{f^{-1}(V)})
\xrightarrow{\theta_V^{-1}}
\mathcal{F}_V
\end{equation}
is equal to $\theta^U_V$. Moreover $(X, \mathcal{F})$ is unique
up to unique isomorphism over $S$.
\end{lemma}

\begin{proof}
By Lemma \ref{lemma-relative-glueing} we get the scheme $X$ over $S$
and the isomorphisms $i_U$.
Set $\mathcal{F}'_U = i_U^*\mathcal{F}_U$ for $U \in \mathcal{B}$.
This is a quasi-coherent $\mathcal{O}_{f^{-1}(U)}$-module.
The maps
$$
\mathcal{F}'_U|_{f^{-1}(V)} =
i_U^*\mathcal{F}_U|_{f^{-1}(V)} =
i_V^*(\rho^U_V)^*\mathcal{F}_U \xrightarrow{i_V^*\theta^U_V}
i_V^*\mathcal{F}_V = \mathcal{F}'_V
$$
define isomorphisms
$(\theta')^U_V : \mathcal{F}'_U|_{f^{-1}(V)} \to \mathcal{F}'_V$
whenever $V \subset U$ are elements of $\mathcal{B}$.
Condition (d) says exactly that this is compatible in case
we have a triple of elements $W \subset V \subset U$ of $\mathcal{B}$.
This allows us to get well defined isomorphisms
$$
\varphi_{12} :
\mathcal{F}'_{U_1}|_{f^{-1}(U_1 \cap U_2)}
\longrightarrow
\mathcal{F}'_{U_2}|_{f^{-1}(U_1 \cap U_2)}
$$
whenever $U_1, U_2 \in \mathcal{B}$ by covering the intersection
$U_1 \cap U_2 = \bigcup V_j$ by elements $V_j$ of $\mathcal{B}$
and taking
$$
\varphi_{12}|_{V_j} =
\left((\theta')^{U_2}_{V_j}\right)^{-1}
\circ
(\theta')^{U_1}_{V_j}.
$$
We omit the verification that these maps do indeed glue to a
$\varphi_{12}$ and we omit the verification of the
cocycle condition of a glueing datum for sheaves
(as in Sheaves, Section \ref{sheaves-section-glueing-sheaves}).
By Sheaves, Lemma \ref{sheaves-lemma-glue-sheaves}
we get our $\mathcal{F}$ on $X$. We omit the verification
of (\ref{equation-glue}).
\end{proof}

\begin{remark}
\label{remark-relative-glueing-functorial}
There is a functoriality property for the constructions explained
in Lemmas \ref{lemma-relative-glueing} and
\ref{lemma-relative-glueing-sheaves}. Namely, suppose given
two collections of data $(f_U : X_U \to U, \rho^U_V)$ and
$(g_U : Y_U \to U, \sigma^U_V)$ as in Lemma \ref{lemma-relative-glueing}.
Suppose for every $U \in \mathcal{B}$ given
a morphism $h_U : X_U \to Y_U$ over $U$ compatible with
the restrictions $\rho^U_V$ and $\sigma^U_V$. Functoriality
means that this gives rise to a morphism of schemes
$h : X \to Y$ over $S$ restricting back to the morphisms $h_U$,
where $f : X \to S$ is obtained from
the datum $(f_U : X_U \to U, \rho^U_V)$ and $g : Y \to S$
is obtained from the datum $(g_U : Y_U \to U, \sigma^U_V)$.

\medskip\noindent
Similarly, suppose given
two collections of data
$(f_U : X_U \to U, \mathcal{F}_U, \rho^U_V, \theta^U_V)$ and
$(g_U : Y_U \to U, \mathcal{G}_U, \sigma^U_V, \eta^U_V)$
as in Lemma \ref{lemma-relative-glueing-sheaves}.
Suppose for every $U \in \mathcal{B}$ given
a morphism $h_U : X_U \to Y_U$ over $U$ compatible with
the restrictions $\rho^U_V$ and $\sigma^U_V$, and a morphism
$\tau_U : h_U^*\mathcal{G}_U \to \mathcal{F}_U$ compatible with
the maps $\theta^U_V$ and $\eta^U_V$. Functoriality
means that these give rise to a morphism of schemes
$h : X \to Y$ over $S$ restricting back to the morphisms $h_U$,
and a morphism $h^*\mathcal{G} \to \mathcal{F}$ restricting back
to the maps $h_U$
where $(f : X \to S, \mathcal{F})$ is obtained from the datum
$(f_U : X_U \to U, \mathcal{F}_U, \rho^U_V, \theta^U_V)$ and
where $(g : Y \to S, \mathcal{G})$ is obtained from the datum
$(g_U : Y_U \to U, \mathcal{G}_U, \sigma^U_V, \eta^U_V)$.

\medskip\noindent
We omit the verifications and we omit a suitable formulation of
``equivalence of categories'' between relative glueing data
and relative objects.
\end{remark}















